\documentclass{article}

% Language setting
% Replace `english' with e.g. `spanish' to change the document language
\usepackage[english]{babel}

% Set page size and margins
% Replace `letterpaper' with `a4paper' for UK/EU standard size
\usepackage[letterpaper,top=2cm,bottom=2cm,left=3cm,right=3cm,marginparwidth=1.75cm]{geometry}

% Useful packages
\usepackage{amsmath, amssymb}
\usepackage{graphicx}
\usepackage[colorlinks=true, allcolors=blue]{hyperref}

\newcommand{\pospart}[1]{\left(#1\right)_{+}}

\begin{document}

\section*{The variational problem}

\subsection*{The primal problem}

With $\beta = 0$, the primal problem is the $\mathrm{min}_{P,Q}$ of:

$$\mathcal{P}(P,Q) := \underbrace{\frac{1}{2}\|K(P)-E\|^2 + \alpha_1\|Q\|_{2,1}}_{F(P,Q)} + \underbrace{\mathcal{I}_0(\nabla_\Omega\mathbb{J}P + \mathrm{div}_Y Q) + {\alpha_2}\|\nabla\mathbb{I}P\|_1}_{G^*(K(P,Q))}$$

where $K$ is the stacked linear operator with matrix representation:

$$\mathcal{K} = \begin{bmatrix} \nabla_\Omega\mathbb{J} & \mathrm{div}_Y \\ \nabla\mathbb{I} & 0 \end{bmatrix} \begin{matrix} \leftarrow f \\ \leftarrow g \end{matrix}$$

$\mathcal{I}_0(\nabla_\Omega\mathbb{J}P + \mathrm{div}_Y Q)$ enforces the Beckmann continuity equation: $Q$ is the displacement-domain flux balancing $\nabla_\Omega \mathbb{J}P$, and $\alpha_1\|Q\|_{2,1}$ is its transport cost. Its multiplier $f$ is the Kantorovich potential. The term ${\alpha_2}\|\nabla\mathbb{I}P\|_1 = \alpha_2\,\mathrm{TV}(\mathbb{I}P)$ regularizes the total mass (with $\alpha_2$ absorbed into $F$ via its conjugate relationship to $g$).

\subsection*{The saddle-point formulation}

In the Fenchel–Rockafellar framework, the problem is written as $\min_{P,Q} F(P,Q) + G^*(\mathcal{K}(P,Q))$, which gives the saddle-point form:

$$\min_{P,Q}\;\max_{f,g}\; F(P,Q) + \langle \mathcal{K}(P,Q),\,(f,g)\rangle - G(f,g)$$

where:

$$G(f,g) = O(f) + \mathcal{I}_{B^\infty}(g/\alpha_2)$$

That is, $F^*$ is finite exactly when $\|g\|_\infty \leq \alpha_2$ and $O(f)=0$ is the constant null function (dual of the indicator in zero $\mathcal{I}_0$).

\subsection*{The dual problem}

From $ F(P,Q) = \tfrac{1}{2} \| K(P)-E \|^2 $ one computes:

$$ F^*(\tilde{P},\tilde{Q}) = \tfrac{1}{2}\| K(K^*K)^{-1}(\tilde{P}+K^*(E))\|^2-\tfrac{1}{2}\|E\|^2+\mathcal{I}_{B^{2\infty}}(\tilde{Q}/\alpha_1) $$

and therefore  since

$$\mathcal{K}^* \begin{pmatrix} f \\ g \end{pmatrix} = -\begin{pmatrix} \mathbb{J}^*\mathrm{div}_\Omega f + \mathbb{I}^*\mathrm{div}\, g \\ \nabla_Y f \end{pmatrix} \begin{matrix} \leftarrow \tilde{P} \\ \leftarrow \tilde{Q} \end{matrix}$$

in the dual problem one will have (by applying $\mathcal{K}^*$ in $G^*(-\mathcal{K}^*(f,g))$, that $\tilde{P}=\mathbb{J}\mathrm{div}_\Omega f + \mathbb{I}^*\mathrm{div} g$ and $\tilde{Q}=\nabla_Yf$. Recall that $\mathbb{J}$ is self-adjoint.

The dual problem is therefore the $\mathrm{min}_{f,g}$ of:

$$ \mathcal{D}(f,g) =
\underbrace{
\tfrac{1}{2}\| K(K^*K)^{-1}(\mathbb{J}\mathrm{div}_\Omega f + \mathbb{I}^*\mathrm{div} g+K^*(E))\|^2 - \tfrac{1}{2}\|E\|^2
+ \mathcal{I}_{B^{2\infty}}(\tfrac{1}{\alpha_1}\nabla_Y f)
}_{F^*(-\mathcal{K}^*(f,g))}
+
\underbrace{
O(f) + \mathcal{I}_{B^\infty}(g/\alpha_2)
}_{G(f,g)}
 $$

---

  This shows why the iterates in iterative primal-dual algorithms evaluated with $\mathcal{P}$ or $\mathcal{D}$ are infinite and therefore it makes *no sense to compute the primal-dual gap* as a stopping criterion.

The gap at iterates $(P^n, Q^n, f^n, g^n)$ is:

$$\mathcal{G}(P^n,Q^n,f^n,g^n) = \underbrace{\mathcal{P}(P^n,Q^n)}_{\text{primal obj. at }(P^n,Q^n)} + \underbrace{\mathcal{D}(f^n,g^n)}_{\text{negative dual obj. at }(f^n,g^n)}$$

Some terms are "proxed", and therefore for example the constraint enforced by $\mathcal{I}_{B^\infty}(g/\alpha_2)$ is satisfied throughout the iterations. While other constraints are not satisfied and will blow up to infinity; that is the case for the terms
$ \mathcal{I}_{B^{2\infty}}(\tfrac{1}{\alpha_1}\nabla_Y f) $ and $ \mathcal{I}_0(\nabla_\Omega\mathbb{J}P + \mathrm{div}_Y Q) $.


\section*{Computing the primal--dual gap via surrogate functions}

\subsection*{The troublesome terms}

Recall from the exact formulation that two indicator terms are infinite at the
iterates of the primal--dual scheme and therefore make the standard gap
unusable as a stopping criterion:
$$
\underbrace{\mathcal{I}_0(\nabla_\Omega\mathbb{J}P + \mathrm{div}_Y Q)}_{\text{primal coupling indicator}}
\qquad\text{and}\qquad
\underbrace{\mathcal{I}_{B^{2\infty}}\!\left(\tfrac{1}{\alpha_1}\nabla_Y f\right)}_{\text{dual indicator from }(\alpha_1\|\cdot\|_{2,1})^{*}} .
$$
The first comes from the coupling constraint $\nabla_\Omega\mathbb{J}P + \mathrm{div}_Y Q = 0$;
its saddle-point partner is $O(f)=\mathcal{I}_0^{*}(f)\equiv 0$ (the multiplier $f$ is unconstrained).
The second is the conjugate of the transport cost $\alpha_1\|Q\|_{2,1}$.
The idea is to replace $\alpha_1\|Q\|_{2,1}$ and the constraint encoded by $\mathcal{I}_0$ by
\emph{surrogate} functions whose conjugates are finite (no indicators), so that both
$\widetilde{\mathcal{P}}$ and $\widetilde{\mathcal{D}}$ are finite at every iterate.
This is licensed by two a~priori bounds, $C_Q$ on $\|Q^{*}\|_{2,1}$ and $C_f$ on $\|f^{*}\|_\infty$.

\subsection*{General picture for the approach}

For a generic problem $\min_x f_1(x_1)+f_2(x_2)+G(\mathcal{K}x)$ with
$\mathcal{K}=[A,\,B]$ and $\mathcal{K}^{*}=\begin{bmatrix}A^{*}\\ B^{*}\end{bmatrix}$,
the dual reads
$$
\min_y\; f_1^{*}(-A^{*}y)+f_2^{*}(-B^{*}y)+G^{*}(y).
$$
If $f_1=\|\cdot\|_\alpha$ (or $f_1=0$), the term $f_1^{*}(-A^{*}y)$ is an indicator
$\mathcal{I}_{B^{\alpha'}}(-A^{*}y)$ that the iterates do not satisfy.
\textbf{Idea:} bound $\|x_1^{*}\|\le C$ in some norm and replace $f_1(x_1)$ by
$f_1(x_1)+\tfrac{1}{2}\pospart{\|x_1\|-C}^{2}$, so that its conjugate becomes regular.
For instance, with $f_1=0$,
$$
\widetilde{f}_1(x_1)=\tfrac{1}{2}\pospart{\|x_1\|-C}^{2}
\qquad\Longrightarrow\qquad
\widetilde{f}_1^{\,*}(-A^{*}y)=\tfrac{1}{2}\big(\|y\|+C\big)^{2}-\tfrac{C^{2}}{2}.
$$

\subsection*{Surrogate primal, dual and saddle-point functionals}

Replacing $\alpha_1\|Q\|_{2,1}$ by $\phi_{C_Q}(\|Q\|_{2,1})$ and
$\mathcal{I}_0(\nabla_\Omega\mathbb{J}P + \mathrm{div}_Y Q)$ by
$C_f\,\|\nabla_\Omega\mathbb{J}P + \mathrm{div}_Y Q\|_1$, the \textbf{primal} becomes
$$
\widetilde{\mathcal{P}}(P,Q)=
\underbrace{\tfrac{1}{2}\|K(P)-E\|^{2} + \phi_{C_Q}(\|Q\|_{2,1})}_{\widetilde{F}(P,Q)}
+ \underbrace{C_f\,\|\nabla_\Omega\mathbb{J}P + \mathrm{div}_Y Q\|_1 + \alpha_2\|\nabla\mathbb{I}P\|_1}_{\widetilde{G}^{*}(\mathcal{K}(P,Q))},
$$
with
$$
\widetilde{F}(P,Q)=\tfrac{1}{2}\|K(P)-E\|^{2} + \phi_{C_Q}(\|Q\|_{2,1}),
\qquad
\widetilde{G}^{*}(a,b)=C_f\|a\|_1 + \alpha_2\|b\|_1 .
$$
Since $(C_f\|\cdot\|_1)^{*}=\mathcal{I}_{C_f B^\infty}$ and $(\alpha_2\|\cdot\|_1)^{*}=\mathcal{I}_{B^\infty}(\cdot/\alpha_2)$,
the conjugate is
$$
\widetilde{G}(f,g)=\mathcal{I}_{C_f B^\infty}(f) + \mathcal{I}_{B^\infty}(g/\alpha_2).
$$
The \textbf{dual} is therefore the $\mathrm{min}_{f,g}$ of
$$
\widetilde{\mathcal{D}}(f,g)=
\underbrace{
\tfrac{1}{2}\big\|K(K^{*}K)^{-1}(\mathbb{J}\mathrm{div}_\Omega f + \mathbb{I}^{*}\mathrm{div}\, g + K^{*}(E))\big\|^{2}
\;-\;\tfrac{1}{2}\|E\|^{2}
\;+\;\phi_{C_Q}^{*}\!\left(\tfrac{1}{\alpha_1}\nabla_Y f\right)
}_{\widetilde{F}^{*}(-\mathcal{K}^{*}(f,g))}
\;+\;
\underbrace{
\mathcal{I}_{C_f B^\infty}(f) + \mathcal{I}_{B^\infty}(g/\alpha_2)
}_{\widetilde{G}(f,g)} ,
$$
where the constant $-\tfrac{1}{2}\|E\|^{2}$ may be dropped for the minimisation.
The associated \textbf{saddle-point} form is
$$
\min_{P,Q}\;\max_{f,g}\;
\tfrac{1}{2}\|K(P)-E\|^{2} + \phi_{C_Q}(\|Q\|_{2,1})
- \mathcal{I}_{C_f B^\infty}(f) - \mathcal{I}_{B^\infty}(g/\alpha_2)
+ \big\langle \mathcal{K}(P,Q),\,(f,g)\big\rangle,
$$
with
$$
\big\langle \mathcal{K}(P,Q),\,(f,g)\big\rangle
= \langle \nabla_\Omega\mathbb{J}P,\,f\rangle + \langle \nabla\mathbb{I}P,\,g\rangle + \langle \mathrm{div}_Y Q,\,f\rangle
= -\langle \mathbb{J}\mathrm{div}_\Omega f,\,P\rangle - \langle \mathbb{I}^{*}\mathrm{div}\, g,\,P\rangle - \langle \nabla_Y f,\,Q\rangle .
$$
The proxed indicators $\mathcal{I}_{C_f B^\infty}(f)$ and $\mathcal{I}_{B^\infty}(g/\alpha_2)$
are now satisfied throughout the iterations (they are enforced by projection),
and the previously divergent terms have been replaced by the finite quantities
$C_f\|\nabla_\Omega\mathbb{J}P + \mathrm{div}_Y Q\|_1$ (primal) and
$\phi_{C_Q}^{*}(\tfrac{1}{\alpha_1}\nabla_Y f)$ (dual).

\subsection*{A priori bounds}

\paragraph{Bound on $Q^{*}$.}
The point $(P^{*},Q^{*})=(0,0)$ is feasible, and by optimality
$$
\tfrac{1}{2}\|K(P^{*})-E\|^{2} + \alpha_1\|Q^{*}\|_{2,1}
+ \mathcal{I}_0(\,\cdot\,) + \alpha_2\|\nabla\mathbb{I}P^{*}\|
\;\le\; \tfrac{1}{2}\|E\|^{2}.
$$
Since every term on the left is nonnegative,
$$
\boxed{\;\|Q^{*}\|_{2,1}\;\le\;\frac{\|E\|^{2}}{2\alpha_1}\;=:\;C_Q\;}.
$$

\paragraph{Bound on $f^{*}$.}
Let the displacement grid be
$$
Y=\Big\{-\tfrac{N_1}{2},\dots,0,\dots,\tfrac{N_1}{2}\Big\}\times\Big\{-\tfrac{N_2}{2},\dots,0,\dots,\tfrac{N_2}{2}\Big\},
\qquad |Y|=(N_1+1)(N_2+1),
$$
with $N_1,N_2$ even (symmetric about $0$). Gauge-fix the potential by
$$
f^{*}(x,0)=0 \qquad \forall x\in\Omega .
$$
Stationarity in $Q$ gives the pointwise gradient bound
$$
\|\nabla_Y f^{*}\|_{2,\infty}\le\alpha_1
\quad\Longrightarrow\quad
\|\nabla f^{*}_x\|_{2,\infty}\le\alpha_1 \quad \forall x\in\Omega,
$$
and summing the increments along a path from the origin to any grid node yields
$$
\boxed{\;\|f^{*}\|_\infty\;\le\;\max\{N_1,N_2\}\Big(1+\tfrac{\sqrt{2}}{2}\Big)\;=:\;C_f\;}.
$$
% NOTE: the increment bound carries a factor alpha_1, so dimensionally one expects
% C_f = alpha_1 * max{N_1,N_2}(1 + sqrt(2)/2). Check whether alpha_1 was set to 1 or absorbed.

\subsection*{The proposed surrogate $\phi_M$ and its convex conjugate}

Take the (reverse-Huber) surrogate of $|\cdot|$
$$
\phi_M(t)=
\begin{cases}
\,|t| & \text{if } |t|\le M,\\[4pt]
\,\dfrac{M}{2}+\dfrac{1}{2M}\,t^{2} & \text{if } |t|> M,
\end{cases}
$$
which coincides with $|t|$ on $[-M,M]$ and grows quadratically outside.

\paragraph{Remark.} $\phi_M$ is even, hence $\phi_M^{*}$ is even; it suffices to take $s\ge 0$.

\paragraph{Conjugate.} By definition
$$
\phi_M^{*}(s)=\sup_t\big(st-\phi_M(t)\big)
=\max\Big\{\;\sup_{|t|\le M}\big(st-|t|\big),\;\;
\sup_{|t|> M}\big(-\tfrac{1}{2M}t^{2}+st-\tfrac{M}{2}\big)\Big\}.
$$
\emph{Branch $|t|\le M$.} For $-1\le s\le 1$ the supremum is $0$ (attained at $t=0$);
the candidate boundary values for $s<-1$ and $s>1$ are dominated below.

\emph{Branch $|t|> M$.} The unconstrained vertex of $-\tfrac{1}{2M}t^{2}+st-\tfrac{M}{2}$ is $t^{*}=sM$:
$$
t^{*}=
\begin{cases}
\,sM & \text{if } |s|\ge 1 \quad(\text{interior}),\\[2pt]
\,M\,\dfrac{s}{|s|} & \text{if } |s|< 1 \quad(\text{at the boundary } t=\pm M),
\end{cases}
\qquad
\text{values}=
\begin{cases}
\,\tfrac{1}{2}M\big(s^{2}-1\big) & |s|\ge 1,\\[4pt]
\,M\big(|s|-1\big) & |s|< 1.
\end{cases}
$$
\emph{Comparison.} For $s>1$: $\tfrac{1}{2}M(s^{2}-1)>M(s-1)\iff \tfrac{1}{2}(s+1)>1\iff s>1$, so the interior value wins;
the symmetric computation holds for $s<-1$. For $|s|\le 1$ the value $0$ dominates $M(|s|-1)\le 0$. Hence
$$
\phi_M^{*}(s)=
\begin{cases}
\,0 & |s|\le 1,\\[4pt]
\,\dfrac{1}{2}M\big(s^{2}-1\big) & |s|\ge 1,
\end{cases}
\qquad\text{or compactly}\qquad
\boxed{\;\phi_M^{*}(s)=\tfrac{1}{2}M\,\pospart{s^{2}-1}\;}
$$
with $\pospart{a}=\max\{a,0\}$ (written $\tfrac{1}{2}M\,|s^{2}-1|_{+}$ in the notes).
Thus $\phi_M^{*}$ is a finite, $C^{1}$ surrogate of $\mathcal{I}_{[-1,1]}$ with ``steepness'' $M$,
and $\phi_M^{*}\to\mathcal{I}_{[-1,1]}$ as $M\to\infty$ --- no indicator appears at finite $M$.




\end{document}
